% ============================================================
%  El protocolo del laboratorio, paso a paso — documento independiente
%  TT 2026-B066 | ESCOM-IPN
%  Frausto Robles Ángel Ali · Rodríguez Verdín Sandoval Vanesa
%
%  Narración cronológica del protocolo de la prueba de nado forzado,
%  desde que se abre el estudio hasta el único dato que el sistema mide.
%  En cada paso se justifica dónde queda registrado ese hecho: primero
%  en el diagrama entidad-relación de la etapa conceptual y después en
%  el grafo relacional de la etapa lógica, ya validado con formas
%  normales.
%
%  Compilar desde la raíz del repositorio:
%      pdflatex -output-directory=build protocolo_fst.tex
% ============================================================

\documentclass[11pt,letterpaper]{article}

\usepackage[spanish]{babel}
\usepackage[T1]{fontenc}
\usepackage[utf8]{inputenc}
\usepackage{lmodern}

\usepackage[
  left=2.4cm, right=2.4cm,
  top=2.2cm,  bottom=2.2cm,
  headheight=15pt,
  headsep=12pt,
  footskip=20pt
]{geometry}

\usepackage{fancyhdr}
\usepackage{lastpage}
\usepackage{xcolor}
\usepackage{array}
\usepackage{tabularx}
\usepackage{enumitem}
\usepackage{titlesec}
\usepackage[hidelinks]{hyperref}

\definecolor{tinta}{HTML}{0D1F21}
\definecolor{suave}{HTML}{3A4E4E}
\definecolor{apagado}{HTML}{677978}
\definecolor{acento}{HTML}{0F6F70}
\definecolor{acentoTenue}{HTML}{DCEBEA}
\definecolor{ambar}{HTML}{A65D2E}
\definecolor{ambarTenue}{HTML}{F4E7DC}
\definecolor{regla}{HTML}{D3DBD8}

\hypersetup{
  pdftitle={El protocolo del laboratorio, paso a paso},
  pdfauthor={Frausto Robles A. A. \& Rodriguez Verdin S. V.},
  pdfsubject={TT 2026-B066 - Protocolo FST contra el modelo de datos}
}

% ---------- Encabezado y pie ----------
\pagestyle{fancy}
\fancyhf{}
\renewcommand{\headrulewidth}{0.4pt}
\renewcommand{\headrule}{\hbox to\headwidth{\color{regla}%
  \leaders\hrule height \headrulewidth\hfill}}
\fancyhead[L]{\footnotesize\sffamily\color{apagado}%
  Protocolo FST · narración cronológica}
\fancyhead[R]{\footnotesize\sffamily\color{apagado}%
  \thepage\ / \pageref{LastPage}}
\fancyfoot[L]{\scriptsize\sffamily\color{apagado}%
  TT 2026-B066 · ESCOM-IPN}
\fancyfoot[R]{\scriptsize\sffamily\color{apagado}%
  Diseño de base de datos · etapas conceptual y lógica}

% ---------- Títulos ----------
\titleformat{\section}
  {\sffamily\bfseries\large\color{tinta}}{}{0pt}{}
  [{\vspace{2pt}\color{tinta}\titlerule[1pt]}]
\titlespacing*{\section}{0pt}{22pt}{10pt}

\setlength{\parindent}{0pt}
\setlength{\parskip}{5pt}

% ---------- Bloques ----------
\newsavebox{\cajaTag}

% \bloque{color}{etiqueta}{texto}
\newcommand{\bloque}[3]{%
  \par\addvspace{6pt}\noindent
  \sbox{\cajaTag}{%
    \begin{minipage}{\dimexpr\linewidth-16pt\relax}
      {\footnotesize\sffamily\bfseries\color{#1}\MakeUppercase{#2}}\par
      \vspace{3pt}
      {\small\color{suave}#3}
    \end{minipage}}%
  \textcolor{#1}{\rule[-\dp\cajaTag]{2pt}{\dimexpr\ht\cajaTag+\dp\cajaTag\relax}}%
  \hspace{12pt}\usebox{\cajaTag}\par\addvspace{6pt}}

\newcommand{\conceptual}[1]{\bloque{acento}{En el esquema conceptual}{#1}}
\newcommand{\relacional}[1]{\bloque{ambar}{En el grafo relacional}{#1}}

% \aviso{etiqueta}{texto}
\newcommand{\aviso}[2]{%
  \par\addvspace{8pt}\noindent
  \fcolorbox{ambar}{ambarTenue}{%
    \begin{minipage}{\dimexpr\linewidth-16pt\relax}
      \vspace{2pt}
      {\footnotesize\sffamily\bfseries\color{ambar}\MakeUppercase{#1}}\par
      \vspace{3pt}
      {\small\color{tinta}#2}
      \vspace{3pt}
    \end{minipage}}\par\addvspace{8pt}}

% \paso{numero}{titulo}
\newcommand{\paso}[2]{%
  \par\addvspace{14pt}\noindent
  {\sffamily\bfseries\color{acento}#1}\hspace{10pt}%
  {\sffamily\bfseries\color{tinta}#2}\par\addvspace{3pt}}

\newcommand{\cod}[1]{{\small\texttt{#1}}}

% ============================================================
\begin{document}

% ---------- Portada breve ----------
{\footnotesize\sffamily\color{acento}%
 TT 2026-B066 · ESCOM-IPN · DISEÑO DE BASE DE DATOS}

\vspace{6pt}
{\LARGE\sffamily\bfseries\color{tinta}%
 El protocolo del laboratorio, paso a paso}

\vspace{4pt}
{\color{tinta}\rule{\linewidth}{1.2pt}}

\vspace{8pt}
{\color{suave}Cada paso dice qué ocurre en el Laboratorio de Bioquímica
Estructural y, enseguida, dónde queda registrado ese hecho: primero en el
diagrama entidad-relación de la etapa conceptual, después en el grafo
relacional de la etapa lógica, ya validado con formas normales.}

\section*{El protocolo, de principio a fin}

% ---------- 0 ----------
\paso{0}{Alguien abre el estudio}

Un investigador decide correr una prueba de nado forzado para evaluar si un
compuesto tiene efecto tipo antidepresivo. Eso es \textbf{un experimento}: un
estudio independiente completo, no una grabación. Según la entrevista (P19),
trae de 4 a 6 grupos, entre 32 y 48 especímenes, y genera alrededor de 8 videos
del Día 2.

\conceptual{\cod{Usuario} $\to$ \emph{registra} (0,N / 1,1) $\to$ \cod{Experimento}.
El (1,1) del lado del experimento dice que \textbf{todo experimento tiene siempre
un creador}; el (0,N) del lado del usuario dice que un investigador puede no
haber creado ninguno.}

\relacional{\cod{EXPERIMENTO(idExperimento, nombre, fecha, notas, idBoleta*)}.
La foránea \cod{idBoleta} es obligatoria: así se implementa el (1,1). Es autoría
y trazabilidad, \textbf{no} control de acceso: cualquier investigador con cuenta
activa ve todos los experimentos del laboratorio (P18, RN-08).}

% ---------- 1 ----------
\paso{1}{Las ratas se reparten en grupos de tratamiento}

El investigador separa a los especímenes en grupos según \textbf{qué recibe cada
uno}: control (sin fármaco), referencia (fluoxetina, el antidepresivo de
referencia) y tratamiento experimental (el compuesto a probar). Mínimo tres
grupos; en la práctica de 4 a 6. Entre 6 y 8 ratas por grupo (PA). Este es el
\textbf{hecho científico central del protocolo}: a qué tratamiento pertenece cada
rata. Todo lo demás es logística.

\conceptual{\cod{Experimento} $\to$ \emph{compone} (3,N / 1,1) $\to$ \cod{Grupo} y
\cod{Grupo} $\to$ \emph{agrupa} (6,8 / 1,1) $\to$ \cod{Espécimen}. El (3,N)
codifica el mínimo de tres grupos; el (6,8) es literalmente el rango que declaró
el laboratorio.}

\relacional{\cod{GRUPO(idGrupo, etiqueta, tipo, tratamiento, idExperimento*)}.
El atributo \cod{tipo} tiene dominio cerrado \{control, referencia, tratamiento
experimental\}; \cod{tratamiento} admite nulo solo si \cod{tipo = control}. La
\cod{etiqueta} distingue grupos del mismo tipo dentro de un experimento, de ahí
la restricción \cod{UNIQUE (idExperimento, etiqueta)}.}

% ---------- 2 ----------
\paso{2}{El grupo no cabe en una sola grabación}

Aquí aparece el nivel que \textbf{ninguna fuente había nombrado}. La cámara web
encuadra máximo 4 cilindros a la vez, y siempre más de 2. Un grupo de 6 a 8 ratas
no cabe. Entonces el laboratorio graba el grupo en \textbf{dos tandas}: primero
3 o 4 ratas, después las otras 3 o 4. La aritmética cierra por tres lados contra
lo que el propio doctor declaró:

\vspace{4pt}
{\small
\begin{tabularx}{\linewidth}{@{}p{4.6cm}p{3.6cm}X@{}}
\multicolumn{3}{@{}l}{\color{regla}\rule{\linewidth}{0.6pt}}\\[-6pt]
{\sffamily\footnotesize\color{apagado}COMPROBACIÓN} &
{\sffamily\footnotesize\color{apagado}RESULTADO} &
{\sffamily\footnotesize\color{apagado}CONTRA LA FUENTE}\\[2pt]
\multicolumn{3}{@{}l}{\color{regla}\rule{\linewidth}{0.6pt}}\\[-4pt]
4 grupos $\times$ 8 ratas & 32 especímenes & P19: «32 a 48»\\
8 ratas $\div$ 4 por cuadro & 2 tandas por grupo & Equipo: máximo 4 a cuadro\\
4 grupos $\times$ 2 tandas & 8 videos de Día 2 & P19: «alrededor de 8»\\[2pt]
\multicolumn{3}{@{}l}{\color{regla}\rule{\linewidth}{0.6pt}}\\
\end{tabularx}}

\conceptual{\cod{Grupo} $\to$ \emph{se graba en} (2,N / 1,1) $\to$ \cod{Tanda} y
\cod{Tanda} $\to$ \emph{aloja} (3,4 / 1,1) $\to$ \cod{Espécimen}. El (2,N) sale de
la división; el (3,4) es el encuadre.}

\relacional{\cod{TANDA(idTanda, ordinal, nCilindros, idGrupo*)}, con
\cod{CHECK nCilindros BETWEEN 3 AND 4} y \cod{UNIQUE (idGrupo, ordinal)}.}

\aviso{Cuidado con el nombre}{«Tanda» es vocabulario nuestro, no del laboratorio.
Es exactamente la pregunta \textbf{P-05} de la reunión de verificación.}

% ---------- 3 ----------
\paso{3}{Cada rata queda identificada y asignada a un cilindro}

El laboratorio identifica físicamente a cada rata (respuesta a \textbf{P-01}).
Además, al montar la grabación, cada rata ocupa un cilindro del encuadre,
contados de izquierda a derecha. Esa posición \textbf{no es lo que identifica a
la rata} —antes sí lo era, y por eso el modelo viejo tenía una cascada de seis
entidades débiles—. Hoy es solo la pista que el pipeline de visión por
computadora usa para saber qué mancha del video corresponde a qué espécimen.

\conceptual{\cod{Espécimen} es \textbf{entidad fuerte}, con \cod{idEspecimen}
propio e \cod{idLaboratorio} como clave alterna. Cuelga de \textbf{dos padres}:
del \cod{Grupo} (qué tratamiento recibió) y de la \cod{Tanda} (en qué grabación
salió).}

\relacional{\cod{ESPECIMEN(idEspecimen, idLaboratorio, numeroCilindro, idTanda*)}.
La foránea \cod{idGrupo} \textbf{no está}: la validación de 3FN encontró que
\cod{idEspecimen} $\to$ \cod{idTanda} $\to$ \cod{idGrupo} era una dependencia
transitiva. Como una tanda pertenece a un solo grupo, saber la tanda ya dice el
grupo, y guardarlo dos veces permitiría que las copias se contradijeran. El grupo
se alcanza uniendo \cod{ESPECIMEN} $\to$ \cod{TANDA} $\to$ \cod{GRUPO}.
Restricciones: \cod{UNIQUE (idTanda, numeroCilindro)}, \cod{UNIQUE idLaboratorio}
(alcance por confirmar, P-03) y \cod{numeroCilindro} $\leq$ \cod{TANDA.nCilindros}.}

\aviso{De qué depende esta corrección}{Toda la eliminación de \cod{idGrupo}
descansa en \textbf{P-04}: si una grabación pudiera mezclar ratas de dos grupos,
la dependencia transitiva se rompe y la foránea \cod{idGrupo} tiene que volver a
\cod{ESPECIMEN}.}

% ---------- 4 ----------
\paso{4}{Día 1 — el estrés agudo}

Se pone a las ratas de una tanda en sus cilindros y se graba una sesión larga, de
unos 20 minutos. Es el pretest: induce el estado que se va a medir al día
siguiente. De este día \textbf{solo se analiza el grupo control} (fuente PC): los
demás se graban, pero su Día 1 no entra al pipeline.

% ---------- 5 ----------
\paso{5}{Día 2 — la medición real, 24 horas después}

Las mismas ratas, la misma tanda, el mismo encuadre, los mismos cilindros. Ahora
se graban \textbf{5 minutos}, y esta es la sesión que produce el dato del
experimento. Que la cámara no se mueva entre los dos días es una garantía
explícita del laboratorio, y es lo que permite comparar Día 1 contra Día 2 del
mismo espécimen.

\conceptual{\cod{Tanda} $\to$ \emph{produce} (1,2 / 1,1) $\to$ \cod{Video}, con
\cod{sesión} en \{Día 1, Día 2\}. El \textbf{(1,2)} es exactamente esto: una tanda
tiene \emph{siempre} su video de Día 2, y \emph{a veces} además el de Día 1.}

\relacional{\cod{VIDEO(idVideo, sesion, archivo, duracion, fechaCarga, idTanda*)},
con \cod{UNIQUE (idTanda, sesion)}: una tanda no puede tener dos videos del mismo
día. \textbf{Segunda divergencia con el diagrama}: en el conceptual la duración
era «20 min Día 1 / 5 min Día 2»; la 3FN detectó que eso era la transitiva
\cod{sesion} $\to$ \cod{duracion}, así que en el esquema definitivo \cod{duracion}
es la duración \emph{real} del archivo y los 20/5 minutos se vuelven verificación
con tolerancia.}

% ---------- 6 ----------
\paso{6}{El video se sube y se procesa}

El investigador carga el MP4. El sistema encola un análisis: preprocesamiento,
detección de cilindros, seguimiento de cada espécimen y clasificación de conducta
por cuadro (visión por computadora y aprendizaje supervisado). Si la confianza de
detección baja de 0.70, el pipeline se detiene en vez de entregar un resultado
silenciosamente incorrecto: por eso importa que la iluminación no esté
garantizada. Un video \textbf{se puede volver a analizar}.

\conceptual{\cod{Video} $\to$ \emph{procesa} (1,N / 1,1) $\to$ \cod{Análisis}.
El (1,N) es la respuesta directa del equipo: «sí, se puede volver a hacer».}

\relacional{\cod{ANALISIS(idAnalisis, estado, etapaActiva, confianza,
nivelClasificacion, fechaAnalisis, idVideo*)}. El atributo
\cod{nivelClasificacion} registra si el clasificador logró separar las tres
conductas o si operó degradado con dos.}

% ---------- 7 ----------
\paso{7}{Se empareja cada rata con el video donde aparece}

Antes de poder decir «esta rata nadó tantos segundos», hay que fijar de qué rata
y de qué video se habla. Ese par —un espécimen concreto en un video concreto— es
una \textbf{observación}. Aquí es donde las dos ramas que se abrieron en el Grupo,
tratamiento por un lado y grabación por el otro, \textbf{se vuelven a juntar}.

\conceptual{\cod{Observación} es \textbf{entidad asociativa} con
\cod{idObservacion} propio y dos interrelaciones: \cod{Espécimen} $\to$
\emph{aparece en} (1,2 / 1,1) y \cod{Video} $\to$ \emph{contiene} (3,4 / 1,1).
El (1,2): una rata aparece en 1 o 2 videos, su Día 2 y quizá su Día 1. El (3,4):
un video contiene una observación por cilindro ocupado. \textbf{Por qué no una
ternaria} Espécimen–Video–Conducta: permitiría emparejar una rata con el video de
otra tanda. La entidad asociativa cierra primero el par válido y clasifica
después.}

\relacional{\cod{OBSERVACION(idObservacion, idEspecimen*, idVideo*)} con
\cod{UNIQUE (idEspecimen, idVideo)}. \textbf{Lo que se perdió al pasar a llaves
subrogadas}: con llaves naturales, las dos foráneas compartían sus primeras
columnas y eso garantizaba \emph{por estructura} que rata y video fueran de la
misma tanda. Ahora hay que declararlo con un \cod{CHECK} o disparador:
\cod{ESPECIMEN.idTanda = VIDEO.idTanda}.}

% ---------- 8 ----------
\paso{8}{Los 5 minutos se parten en minutos}

El laboratorio no quiere un total: quiere el desglose minuto a minuto, porque la
curva importa (P9). Los 5 minutos del Día 2 dan exactamente cinco tramos.

\conceptual{\cod{Observación} $\to$ \emph{se divide en} (5,5 / 1,1) $\to$
\cod{Intervalo}, con \cod{minuto} en 1..5. El (5,5) es una cardinalidad
\textbf{fija}, no un rango: siempre cinco.}

\relacional{\cod{INTERVALO(idIntervalo, minuto, idObservacion*)} con
\cod{UNIQUE (idObservacion, minuto)} y \cod{CHECK minuto BETWEEN 1 AND 5}.}

% ---------- 9 ----------
\paso{9}{En cada minuto se contabilizan los segundos por conducta}

Y aquí, por fin, \textbf{el único dato que el sistema mide en todo el protocolo}:
cuántos segundos pasó esa rata, en ese minuto, nadando, inmóvil o escalando. El
buceo cuenta como nado activo (P7); un episodio cuenta si dura de 3 a 5 segundos
continuos (P6). Las tres conductas son mutuamente excluyentes por cuadro.

\conceptual{\cod{Intervalo} $\to$ \emph{presenta} (2,3 / 0,N) $\to$ \cod{Conducta},
con \cod{segundos} como \textbf{atributo de la interrelación}. El (2,3) no es
capricho: en modo normal son tres conductas, en modo degradado el catálogo baja
a dos.}

\relacional{\cod{PRESENTA(idIntervalo*, conducta*, segundos)}, llave primaria de
\textbf{dos} componentes. En el modelo viejo esta misma tabla terminaba con una
llave de \textbf{siete} columnas y una sola de dato; ese fue el argumento
decisivo para adoptar llaves subrogadas.}

% ---------- 10 ----------
\paso{10}{Lo que el investigador consulta (y que no se guarda)}

Total por conducta y por espécimen, comparación Día 1 contra Día 2, media,
desviación estándar y varianza por grupo, comparación entre grupos, exportación
a CSV, XLSX y PDF.

\relacional{\textbf{Nada de eso está almacenado}: todo se deriva de
\cod{PRESENTA.segundos} sumando hacia arriba por el grafo. Es la actividad 6 del
método conceptual, sin redundancia. Si el desempeño lo exigiera, materializarlos
es una decisión de la \textbf{etapa física}, donde el método sí admite
«redundancia controlada», y justificada por medición, no antes.}

% ============================================================
\section*{Dónde el grafo relacional no se lee igual que el diagrama}

Son las tres correcciones que encontró la validación con formas normales.
Conviene tenerlas claras antes de la revisión de notación, porque son justo lo
que puede preguntar la Dra. Cordero.

\vspace{4pt}
{\small
\begin{tabularx}{\linewidth}{@{}p{4.4cm}p{4.2cm}X@{}}
\multicolumn{3}{@{}l}{\color{regla}\rule{\linewidth}{0.6pt}}\\[-6pt]
{\sffamily\footnotesize\color{apagado}EN EL DIAGRAMA E-R} &
{\sffamily\footnotesize\color{apagado}EN EL GRAFO RELACIONAL} &
{\sffamily\footnotesize\color{apagado}POR QUÉ}\\[2pt]
\multicolumn{3}{@{}l}{\color{regla}\rule{\linewidth}{0.6pt}}\\[-4pt]
\cod{Espécimen} tiene dos padres: Grupo y Tanda &
\cod{ESPECIMEN} tiene una sola foránea: \cod{idTanda} &
3FN: \cod{idGrupo} era transitiva. El grupo se alcanza por la tanda\\[4pt]
\cod{Video.duración} = 20 min / 5 min según sesión &
\cod{VIDEO.duracion} = duración real del archivo &
3FN: \cod{sesion} $\to$ \cod{duracion} era transitiva\\[4pt]
\cod{Usuario.nombre} &
\cod{USUARIO.nombre} + \cod{USUARIO.apellidos} &
1FN: era un dominio compuesto\\[2pt]
\multicolumn{3}{@{}l}{\color{regla}\rule{\linewidth}{0.6pt}}\\
\end{tabularx}}

\textbf{La semántica del protocolo no cambia en ninguna de las tres.} La rata
sigue perteneciendo a un tratamiento; simplemente el esquema no guarda ese hecho
dos veces.

% ============================================================
\section*{Una pieza del protocolo que el esquema todavía no representa}

\cod{OBSERVACION} cuelga de \cod{VIDEO}, pero la cardinalidad \cod{VIDEO} $\to$
\cod{ANALISIS} es (1,N). Si se reanaliza un video y el resultado cambia, ese
segundo resultado \textbf{no tiene dónde escribirse}, y desde una fila de
\cod{PRESENTA} no se puede saber qué análisis la produjo ni con qué nivel de
clasificación.

\bloque{acento}{Corrección propuesta}{Que \cod{OBSERVACION} referencie
\cod{idAnalisis} en lugar de \cod{idVideo}; el video se alcanza por
\cod{ANALISIS} $\to$ \cod{VIDEO}. No cambia el número de relaciones ni rompe BCNF.}

\aviso{No aplicar todavía}{Depende de \textbf{P-08}. Si el laboratorio prefiere
que el reanálisis reemplace al anterior, el esquema se queda como está, pero se
pierde la evidencia de repetibilidad, que es justo el criterio que el laboratorio
pidió como principal.}

\end{document}
